\subsection{Conclusion and outlook}

The precision requirements of modern $\ell p$ scattering experiments require unprecedented theoretical calculations.
A critical part of this is the development of event generators that can be used in experimental simulations.
These requirements mean that we need to go beyond leading order for the radiative process $\ell p\to \ell p\gamma$. This is also roughly equivalent to going beyond NLO for $\ell p\to \ell p$.
Currently, no event generator with these capabilities exist even though some codes (notably {\sc McMule} and {\sc MESMER}) could be adapted.

Unfortunately, the increased precision comes with a massive increase of the complexity of the distributions that need to be sampled.
Luckily this is a known problem;
the LHC community has worked on issues that are almost identical for decades.
While not all ideas are immediately transferable to our process, many are.
Additionally to the calculation of one-loop matrix elements (which is fully automatised for all standard theories), this includes methods to deal with IR singularities, negative weights, and numerical resummation using a parton shower.

Using these tried and tested techniques will allow us to construct event generators with (and beyond) NNLO accuracy for $\ell p\to \ell p$ in the coming years.
